% ============================================================================
% บทที่ 8 — บันทึกและแชร์บนคลาวด์
% ============================================================================
\mychapter{บันทึกและแชร์บนคลาวด์}

\noindent บทนี้อธิบายการบันทึกโครงการขึ้นคลาวด์ การสร้างลิงก์แชร์ การจัดการสิทธิ์
การรับชมจากผู้รับ และการจัดการข้อขัดแย้งของฉบับ

ฟังก์ชันบันทึกและแชร์จะใช้ได้เมื่อระบบเปิดใช้งานบริการคลาวด์แล้ว

\begin{caution}
หากยังไม่มีการตั้งค่าบริการคลาวด์ ฟังก์ชันบันทึกและแชร์จะไม่ทำงาน
ควรติดต่อผู้ดูแลระบบเพื่อขอให้ตั้งค่าบริการก่อน
\end{caution}

\section{การบันทึกโครงการขึ้นคลาวด์}

\begin{steps}
  \item ไปที่หน้า \textbf{บันทึกและแชร์}
  \item การบันทึกครั้งแรก ระบุ \textbf{ชื่อโครงการ} เช่น
        \texttt{โครงการตารางสอบตัวอย่าง} แล้วกดปุ่ม \texttt{บันทึกขึ้นคลาวด์}
  \item การบันทึกครั้งต่อไป กดปุ่ม \texttt{บันทึกเดี๋ยวนี้}
\end{steps}

\begin{expected}
ระบบสร้างโครงการบนคลาวด์ และทุกครั้งที่บันทึกจะสร้างฉบับใหม่ของโครงการ
หน้าจอแสดงชื่อโครงการ ฉบับปัจจุบัน และเวลาที่บันทึกล่าสุด
การบันทึกครั้งต่อไปสามารถทำได้ทุกเมื่อ ไม่ต้องรอให้แก้ไขถึงจุดใดจุดหนึ่งก่อน
\end{expected}

\manualfigure{docs/usage-report/screenshots/crops/F26b.png}{0.95\textwidth}{%
  หน้าบันทึกและแชร์ แสดงโครงการบนคลาวด์และเวลาที่บันทึกล่าสุด}

เมื่อเปิดใช้งานตัวเลือก \textbf{บันทึกอัตโนมัติเมื่อมีการแก้ไข}
ระบบจะบันทึกฉบับใหม่หลังจากหยุดแก้ไขชั่วครู่ ไม่ต้องกดบันทึกเอง

\begin{note}
การบันทึกอัตโนมัติทำงานเฉพาะเมื่อเปิดตัวเลือกไว้
เมื่อจะปิดหน้าจอหรือหยุดงาน ควรตรวจว่าการแก้ไขล่าสุดถูกบันทึกแล้ว
\end{note}

\section{การสร้างลิงก์แชร์}

\begin{caution}
ถือว่าลิงก์แชร์เป็นข้อมูลลับ โดยเฉพาะลิงก์สิทธิ์ \texttt{โครงการเต็ม} และ
\texttt{แก้ไข} ควรเลือกสิทธิ์ต่ำสุดที่เพียงพอ กำหนดวันหมดอายุเมื่อเหมาะสม
และไม่ส่งลิงก์ดังกล่าวในพื้นที่สาธารณะ
\end{caution}

\begin{steps}
  \item ไปที่หน้า \textbf{บันทึกและแชร์}
  \item ในส่วน \textbf{ลิงก์แชร์} ให้กดปุ่ม \texttt{สร้างลิงก์แชร์}
  \item เลือก \textbf{สิทธิ์} ที่ต้องการ
  \item เลือก \textbf{ฉบับ} (\texttt{ตามฉบับล่าสุด} หรือ \texttt{ตรึงฉบับที่บันทึกไว้})
  \item เลือก \textbf{อายุ} (\texttt{ไม่มีวันหมดอายุ} หรือ 1/7/30 วัน)
  \item กดปุ่ม \texttt{สร้างลิงก์แชร์}
\end{steps}

\begin{expected}
ระบบสร้างลิงก์แชร์และแสดงในรายการ พร้อมปุ่ม \texttt{คัดลอกลิงก์}
เพื่อคัดลอกไปส่งให้ผู้รับ และปุ่ม \texttt{เพิกถอน}
\end{expected}

\begin{note}
การสร้างลิงก์แชร์จะส่งข้อมูลจากฉบับที่บันทึกไว้บนคลาวด์เท่านั้น
ไม่รวมการแก้ไขที่ยังไม่ได้บันทึกในเครื่อง
\end{note}

\section{สิทธิ์ของลิงก์แชร์}

ลิงก์แชร์มีสี่ระดับสิทธิ์:

\begin{center}
\small
\begin{tabular}{L{3.2cm}cccc}
\toprule
\textbf{สิทธิ์} & \textbf{ดูตาราง} & \textbf{ตรวจสอบ} & \textbf{ข้อมูลโครงการ} & \textbf{จัดการแชร์} \\
\midrule
\texttt{ตารางสอบ} & \yesmark & \nomark & \nomark & \nomark \\
\texttt{ตรวจสอบ} & \yesmark & \yesmark & \nomark & \nomark \\
\texttt{โครงการเต็ม} & \yesmark & \yesmark & \yesmark & \nomark \\
\texttt{แก้ไข} & \yesmark & \yesmark & \yesmark & \yesmark \\
\bottomrule
\end{tabular}
\end{center}

\begin{description}
  \item[\texttt{ตารางสอบ}] ดูตารางสอบได้เท่านั้น
  \item[\texttt{ตรวจสอบ}] ดูตารางสอบและข้อมูลการตรวจสอบได้
  \item[\texttt{โครงการเต็ม}] ดูข้อมูลทั้งหมด รวมถึงไฟล์ต้นทาง (ถ้ามี)
        แต่ไม่สามารถแก้ไขได้
  \item[\texttt{แก้ไข}] แก้ไขข้อมูลและจัดการลิงก์แชร์ได้
\end{description}

\section{ฉบับของลิงก์แชร์}

เมื่อสร้างลิงก์แชร์ สามารถเลือกฉบับที่ลิงก์อ้างอิงได้:
\texttt{ตามฉบับล่าสุด} จะแสดงฉบับล่าสุดเสมอเมื่อผู้รับเปิดลิงก์
ส่วน \texttt{ตรึงฉบับที่บันทึกไว้} จะแสดงฉบับตอนสร้างลิงก์เท่านั้น
แม้จะมีฉบับใหม่กว่าในภายหลัง

\section{การเพิกถอนลิงก์}

\begin{steps}
  \item ในรายการลิงก์แชร์ ให้กดปุ่ม \texttt{เพิกถอน} ถัดจากลิงก์ที่ต้องการ
\end{steps}

\begin{caution}
การเพิกถอนลิงก์จะทำให้เปิดโครงการต่อจากลิงก์นั้นไม่ได้อีก
แต่เป็นการควบคุมการเข้าถึงในอนาคตเท่านั้น
ไม่ได้ลบไฟล์ที่ผู้รับดาวน์โหลดไปแล้ว
\end{caution}

\section{การเปิดลิงก์แชร์จากผู้รับ}

\begin{steps}
  \item ผู้รับเปิดลิงก์แชร์ที่ได้รับในเบราว์เซอร์
  \item อ่านแถบแสดงสิทธิ์ที่ด้านบนของหน้า ว่าตัวเองดูอะไรได้บ้าง
\end{steps}

\begin{expected}
ลิงก์ \texttt{ตารางสอบ} หรือ \texttt{ตรวจสอบ} จะแสดงแบบอ่านอย่างเดียว
ลิงก์ \texttt{แก้ไข} จะแสดงปุ่ม \texttt{แก้ไขต่อในเครื่องนี้}
เมื่อกดแล้วผู้รับสามารถแก้ไขข้อมูลและบันทึกฉบับใหม่ได้
\end{expected}

\manualfigure{docs/usage-report/screenshots/crops/F27b.png}{0.95\textwidth}{%
  แถบด้านบนของหน้าที่เปิดจากลิงก์แชร์แบบอ่านอย่างเดียว:
  ระบุว่า \texttt{ตารางสอบที่แชร์} และ \texttt{อ่านอย่างเดียว}}

\manualfigure{docs/usage-report/screenshots/crops/F28b.png}{0.95\textwidth}{%
  แถบด้านบนของหน้าที่เปิดจากลิงก์แชร์แบบแก้ไข:
  ระบุว่า \texttt{สิทธิ์แก้ไขที่แชร์}}

\begin{note}
ลิงก์ที่เลือก \texttt{ตามฉบับล่าสุด} จะเห็นการแก้ไขล่าสุดก็ต่อเมื่อ
เจ้าของโครงการบันทึกฉบับขึ้นคลาวด์แล้ว และผู้รับโหลดลิงก์ใหม่
\end{note}

\needspace{8\baselineskip}
\section{การจัดการข้อขัดแย้งของฉบับ}

เมื่อมีสองอุปกรณ์แก้ไขโครงการเดียวกัน และอุปกรณ์หนึ่งบันทึกขึ้นคลาวด์ก่อน
อุปกรณ์อื่นที่ยังถือฉบับเก่าอยู่จะพบหน้าต่างข้อขัดแย้งเมื่อพยายามบันทึก

\begin{expected}
ระบบแสดงหน้าต่าง \textbf{มีฉบับอื่นถูกบันทึกก่อน} พร้อมข้อความว่า
\textbf{อุปกรณ์อื่นบันทึกโครงการนี้ขึ้นคลาวด์ก่อนเบราว์เซอร์นี้
ระบบจะไม่เขียนทับเงียบ ๆ โปรดเลือก}
\end{expected}

\manualfigure{docs/usage-report/screenshots/crops/F30b.png}{0.72\textwidth}{%
  หน้าต่างข้อขัดแย้งฉบับ \textbf{มีฉบับอื่นถูกบันทึกก่อน}}

ระบบจะแสดงตัวเลือกสามแบบ:

\begin{description}
  \item[\textbf{ดาวน์โหลด project.zip}] เก็บสำเนาการแก้ไขในเครื่องไว้เป็นไฟล์
        เพื่อตรวจสอบหรืออ้างอิงภายหลัง (ไฟล์นี้นำกลับเข้าระบบไม่ได้ ดูบทที่ 7)
  \item[\textbf{แก้ต่อในเครื่องนี้}] เก็บการแก้ไขในเครื่องไว้ต่อ
        โดยไม่บันทึกขึ้นคลาวด์ (ไม่ได้แก้ข้อขัดแย้ง หรือรวมการแก้ไขเข้าด้วยกัน)
  \item[\textbf{โหลดฉบับจากคลาวด์}] โหลดฉบับจากคลาวด์มาแทนที่
        การแก้ไขในเครื่องที่ยังไม่ได้บันทึก
\end{description}

\begin{caution}
การกด \texttt{โหลดฉบับจากคลาวด์} จะแทนที่การแก้ไขในเครื่องที่ยังไม่ได้บันทึก
หากยังต้องการใช้การแก้ไขนั้น ควรกด \texttt{ดาวน์โหลด project.zip} เก็บไว้ก่อน
\end{caution}

\seealso{บทที่ 7 สำหรับไฟล์เก็บถาวรและการเปิดงานต่อ และบทที่ 5 สำหรับการแก้ไขเวลา}
